\section{Related work}

The closest literature tests whether conditional average treatment effects are constant across covariate-defined subpopulations. \citet[Section III]{CrumpHotzImbensMitnik2008Homogeneity} develop a sieve Wald approach that compares series estimates of the two conditional outcome means and removes the constant component when testing homogeneity. Their analysis combines smoothness and sieve-growth restrictions with conditions on conditional variances and third moments of outcome disturbances to obtain asymptotic calibration and consistency against specified smooth alternatives. This establishes an important precedent for treating homogeneity as a global restriction on conditional means. The present paper studies that restriction through finite-sample control over the entire composite null and a population separation criterion that is uniform over the stipulated model class; the corresponding guarantees are stated in \cref{obj:thm:whole-null-calibration-resolved,obj:thm:original-record-attainable-rate}.

Orthogonal conditional moments provide a second route to the same substantive question. \citet[Assumptions 1--5, Proposition 2, and Theorems 3.1--3.3 and 4.1]{LuSong2026OrthogonalICM} construct a marked empirical process from a doubly robust score and centered covariate indicators. Their theory specifies series estimators for the propensity score and outcome regressions, with primitive smoothness and series-growth conditions ensuring a feasible-to-oracle approximation uniform over the process index. The outcome conditions include finite second moments and uniformly bounded conditional variances. The resulting tests have asymptotic Gaussian-process calibration, multiplier-bootstrap validity, consistency against fixed alternatives, and nontrivial power along specified local alternatives approaching the null at the inverse square root of sample size. Uniformity over the process index supports these distributional approximations. Our separation criterion instead takes a supremum over all admissible population laws whose effect function is sufficiently far from its average, as specified in \cref{obj:def:testing-risk}.

\citet[Assumptions 3.3--3.4, Proposition 3.1, and Section 3.3]{LapentaStrittmatterVergara2026Omnibus} integrate orthogonal unconditional moments to obtain an omnibus test that accommodates flexible nuisance estimation. Their formulation separates the controls used for identification from the covariates used to assess heterogeneity. Under unconfoundedness, the stated theory assumes bounded outcomes and nuisance estimators converging faster than the inverse fourth root of sample size in population \(L_2\), together with boundedness and overlap conditions on the fitted functions. It establishes an asymptotic null distribution, consistency against fixed alternatives, and a bootstrap implementation that retains the fitted nuisance functions across replications. These conditions make explicit the nuisance-learning and outcome envelopes supporting their calibration. Here, the primitive restrictions in \cref{obj:def:model} determine the public tuning and finite-sample bounds directly.

A closely related scalar target is the variance of the conditional average treatment effect. \citet[Assumptions 3--4, Algorithm 1, and Theorem 1]{Yu2026VarianceHTE} address its zero boundary using intra-fold sample splitting: distinct evaluation subsamples share the same out-of-fold nuisance estimates, producing cancellation of nuisance-bias terms. The stated asymptotic normality result uses finite fourth outcome moments, a nondegeneracy condition, and convergence faster than the inverse fourth root of sample size for the propensity, outcome-regression, and conditional-effect estimators. Under the uniform covariate marginal in \cref{obj:ass:uniform}, the squared population \(L_2\) distance of the effect from its average also equals its variance. Our testing risk therefore measures separation using a closely related population quantity, with calibration and power assessed uniformly over the model's conditional-moment envelope.

Earlier variance-target work studies estimation and confidence intervals rather than minimax separation. \citet[Sections 2.1--2.3]{Levy2021} define the variance of the full conditional treatment-effect function and develop TMLE and cross-validated TMLE estimators; their efficiency conditions include control of an outcome-regression remainder that is not removed merely by knowing the treatment mechanism. \citet[Sections 3--4, Theorems 6--7]{SanchezBecerra2023} instead work in experiments with known assignment probabilities and a predictive first stage. They establish efficiency for the full variance under additional learning conditions and uniform asymptotic coverage for a fold-specific proxy variance under their stated fourth-moment and regularity conditions. Our target is the same full population variance under the uniform design, but the criterion is finite-sample testing risk over Hölder and conditional-moment classes; it does not replace the variance by a fold-specific proxy or assume a learned first stage.

Directional testing offers a complementary way to concentrate power. \citet[Section 2.2.1, Assumptions 1--5, Theorem 1, and Corollaries 1--2]{DhawanGuoShah2026Directional} learn a promising score direction on one sample and evaluate a debiased score on an independent sample. Their general theory incorporates conditional score-moment bounds, variance conditions, uniform nuisance-estimation requirements, and the additional moment conditions stated in their asymptotic theorem. It delivers uniform asymptotic null calibration and a power guarantee requiring sufficient correlation between the learned direction and the population signal. Thus, uniform asymptotic inference is already an explicit objective in this literature. The present analysis adds finite-sample whole-null calibration and uniform separation guarantees indexed by primitive moment and smoothness restrictions. The distinction concerns both the error criterion and the alternatives over which power is guaranteed.

The quadratic structure connects these causal tests to nonparametric goodness-of-fit theory. \citet[Sections 2--4 and Theorems 1--2]{IngsterSapatinas2009GoodnessOfFit} analyze minimax separation in multivariate regression using increasing-dimensional U-statistics and ellipsoidal smoothness classes. \citet[Sections 2.3--2.4 and Theorems 1--2]{CommingesDalalyan2013QuadraticNull} study composite regression nulls defined by quadratic functionals, obtaining sharp asymptotic testing results for nonnegative diagonal functionals under their stated conditions. These analyses explain how approximation error and the stochastic cost of a quadratic statistic jointly determine detectable signal size. In the present setting, that balance also incorporates unknown treatment assignment and conditional outcome tails. The comparisons in \cref{obj:thm:heavy-tail-comparison,obj:thm:exact-propensity-preservation-boundary} characterize their consequences for critical separation radii within the stipulated classes.

The projection corrections also have a close methodological antecedent in higher-order causal estimation. \citet[Sections 3--4]{KennedyBalakrishnanRobinsWasserman2024MinimaxCATE} derive minimax bounds for pointwise conditional-effect estimation in Hölder models and construct a higher-order local polynomial R-learner, with attainment under the stated additional conditions. Its second-order corrections use projections to control nuisance-product bias, linking the accuracy of causal estimation to separate smoothness restrictions on treatment assignment and outcome means. Our corrected scores in \cref{obj:def:score-family,obj:def:explicit-score-ledger} use related projection ideas for a global homogeneity test. The population targets and loss criteria distinguish the two analyses: pointwise estimation measures recovery of the effect function at a covariate value, whereas \cref{obj:def:testing-risk,obj:def:critical-radius} assess rejection errors for constancy against population-separated alternatives.

Outcome-tail handling connects the construction to robust U-statistics. \citet[Section 2 and Theorem 3]{JolyLugosi2016RobustU} use median-of-means aggregation to estimate U-statistic expectations; their canonical-statistic result permits moment orders strictly above one and at most two. This supplies a precise comparator for finite-moment control of nonlinear statistics. Our construction combines clipping, projection corrections, and independent evaluation blocks, with explicit bias and covariance bounds in \cref{obj:def:explicit-score-ledger}. Together with the full-record converse in \cref{obj:thm:tail-essential-full-record-converse}, the analysis relates the attainable separation scale to the conditional outcome-moment restriction.

Other heterogeneity procedures address distinct scientific targets. \citet[Sections 2--4]{DukesStensrudBrioschiHudson2026Homogeneity} study quantitative and qualitative effect heterogeneity through contrasts tied to the population value of treatment rules, with asymptotic testing guarantees under their regularity conditions. \citet[Sections 2--3 and Theorem 3.3]{JainLuedtke2026Distributional} target equality of conditional potential-outcome distributions and establish asymptotic bootstrap calibration and consistency against fixed alternatives under their stated conditions. \citet[Sections 2--3]{AshouriHenderson2026Predictive} assess the predictive advantage of fitted models that allow heterogeneous effects over models imposing a constant effect. Policy-value contrasts, conditional distributional equality, and predictive performance each give heterogeneity a specific operational meaning. The contribution developed here concerns constancy of the conditional mean contrast under fixed overlap, a known uniform covariate marginal, stipulated smoothness restrictions, and a conditional raw-moment envelope. Within that scope, it combines whole-null finite-sample calibration with separation-rate comparisons that quantify the roles of outcome tails and supplied propensity information.
